\appendix

\onecolumn

\section{Vehicle Routing Problem and Its Decomposition} \label{apx:vhr-decomposition}
Consider the vehicle routing problem, where 
\begin{itemize}
    \item $V$ denotes the set of nodes, each of which represents a customer or the depot,  
    \item $K$ is the set of vehicles,
    \item $c_{ij}$ is the cost (e.g., distance) to travel from node $i$ to node $j$, where $i, j \in V$,
    \item $d_i$ denotes the demand of customer $i$, and
    \item $Q$ denotes the capacity of each vehicle (assume that all vehicles $k \in K$ have the same capacity).
\end{itemize}
The decision variables is $x \in \{0, 1\}^{\abs{V} \times \abs{V} \times \abs{K}}$, where $x_{ijk}$ equals to 1 if the vehicle $k \in K$ travels directly from node $i$ to node $j$, and 0 otherwise. Then, the vehicle routing problem is formulated as follows
\begin{subequations}
    \begin{align}
    \min & \displaystyle \sum_{i\in V, j\in V, k \in K} c_{ij}x_{ijk} \notag \\
        \mathrm{s.t.~} & x \in \{0, 1\}^{\abs{V} \times \abs{V} \times \abs{K}} \\
        & \displaystyle \sum_{k \in K} \sum_{j \in V} x_{jik} = 1 & \forall i\in V \setminus \{0\} \label{eq:vrp-linking}\\
        & \displaystyle \sum_{j \in V\setminus\{0\}}x_{0jk} = 1 & \forall k \in K \label{eq:vrp-preservation1}\\
        & \displaystyle \sum_{i \in V \setminus \{0\}} x_{i0k} = 1 &\forall  k \in K \label{eq:vrp-preservation2}\\
        & \displaystyle \sum_{j \in V} x_{jik} - \sum_{j \in V} x_{ijk} = 0 & \forall i\in V \setminus \{0\}, \forall k \in K \label{eq:vrp-flow-conservation}\\
        & \displaystyle \sum_{i \in V \setminus \{0\}} d_i \Big(\sum_{j \in V} x_{jik}\Big) \leq Q & \forall k \in K . \label{eq:vrp-capacity-restriction}
    \end{align}
\end{subequations}
Here, the constraint \eqref{eq:vrp-linking} is the linking constraint, ensuring that each customer is visited exactly once by a single vehicle. The constraints \eqref{eq:vrp-preservation1} and \eqref{eq:vrp-preservation2} ensure that every vehicle must depart from the depot and must return to the depot. The constraint \eqref{eq:vrp-flow-conservation} means that if a vehicle arrives at a customer, it must also depart from the same customer. And the final constraint \eqref{eq:vrp-capacity-restriction} implies that the total customer demand on any single vehicle's route must not exceed its capacity. We have omitted the subtour elimination constraints for each vehicle to simplify exposition. 

Note that the constraint \eqref{eq:vrp-linking} is the crucial link that imposes the relation between vehicles. If we dualize such a constraint with Lagrangian multipliers, we can decompose the original complicated problems into many easier-to-solve sub-problems, each corresponding to a vehicle $k \in K$. Concretely, let $f(\boldsymbol{\pi})$, where $\pi \in \bbR^{\abs{V} - 1}$, be the objective function of the Lagrangian relaxation problem corresponding to the multiplier $\pi$, we have
\begin{equation*}
    \begin{aligned}
        f(\pi) &= \sum_{i \in V, j \in V, k \in K}c_{ij} x_{ijk} + \sum_{i \in V \setminus\{0\}}\pi_i\left(\sum_{k \in K} \sum_{j\in V}x_{jik} - 1\right) \\
        &= \sum_{k \in K}\left(\sum_{i \in V, j \in V} c_{ij} x_{ijk} + \sum_{i\in V\setminus \{0\}, j \in V} \pi_i x_{jik}\right) - \sum_{i \in V \setminus \{0\}} \pi_i.
    \end{aligned}
\end{equation*}
Moreover, note that the other constraints (constraints \eqref{eq:vrp-preservation1}, \eqref{eq:vrp-preservation2}, \eqref{eq:vrp-flow-conservation}, \eqref{eq:vrp-capacity-restriction}) are defined for each vehicle $k \in K$. Therefore, given the multiplier $\pi$, we can decompose the original problems into many sub-problems $\boldsymbol{P}_k$, for $k \in K$, as follow
\begin{equation}
    \begin{array}{cl}
        \min & \displaystyle \sum_{i \in V, j \in V}c_{ij} x_{ijk} + \sum_{i\in V\setminus \{0\}, j \in V}\pi_i x_{jik} \\
        \mathrm{s.t.~} &  x_k \in \{0, 1\}^{\abs{V} \times \abs{V}} \\
        &\displaystyle \sum_{j \in V\setminus\{0\}} x_{0jk} = 1, \quad \sum_{i \in V \setminus \{0\}} x_{i0k} = 1\\
        & \displaystyle \sum_{j \in V} x_{jik} - \sum_{j \in V} x_{ijk} = 0 \quad \forall i\in V \setminus \{0\} \\
        & \displaystyle \sum_{i \in V \setminus \{0\}} d_i \left(\sum_{j \in V} x_{jik}\right) \leq Q.
    \end{array}
    \tag{$\boldsymbol{P}_k$}
\end{equation}

\section{Additional Backgrounds} \label{apx:additional-backgrounds}
First, we recall Talagrand's contraction inequality, which is crucial for analyzing the Rademacher complexity of a class of composite functions involving Lipschitz functions. 
\begin{lemma}[Talagrand's contraction inequality, \citet{ledoux1991probability}] \label{lm:talagrand}
    Let $\cU$ be a real-valued function class that take inputs from domain $\cP$, and let $S = \{P_1, \dots, P_N\} \subset \cP$ be a set of $N$ samples. Let $\phi_1, \dots, \phi_N$ be a set of $L$-Lipschitz functions (i.e., $\abs{\phi_i(a) - \phi_i(b)} \leq L \abs{a - b}$ for all $a, b \in \bbR$). Then we have 
    \[
        \bbE_{\sigma}\left[\sup_{u \in \cU} \frac{1}{N}\sum_{i = 1}^N \sigma_i \phi_i(u(P_i))\right] \leq L \cdot \bbE_{\sigma} \left[\sup_{u \in \cU}\frac{1}{N}\sum_{i = 1}^N \sigma_i u(P_i)\right],
    \]
    where $\sigma = (\sigma_1, \dots, \sigma_N)$ are i.i.d.~Rademacher random variables. 
\end{lemma}

We recall the vector Bernstein's inequality, a concentration inequality with controlled variance. 
\begin{lemma}[Vector Bernstein's inequality, \citet{bernstein1924modification}] \label{lm:vector-bernstein}
    Let $X_1, \dots, X_N$ be independent zero-mean random vectors in $\mathbb{R}^d$. Suppose that $\|X_i\|_2 \le M$ almost surely for all $i \in \{1, \dots, N\}$. Let $\sigma^2 = \sum_{i=1}^N \mathbb{E}[\|X_i\|_2^2]$ be the total variance. Then for any $t > 0$,$$\mathbb{P}\left( \left\| \sum_{i=1}^N X_i \right\|_2 > t \right) \le (d+1) \exp\left( \frac{-t^2}{2(\sigma^2 + Mt/3)} \right).$$
\end{lemma}
We then recall Popoviciu's inequality, a preliminary result that provides an upper bound for the variance of random variables with bounded support. 

\begin{lemma}[Popoviciu's inequality, \citet{popoviciu1965certaines}] \label{lm:popoviciu}
    Let $X$ be a random variable such that it is bounded by a minimum value $m$ and a maximum value $M$. Then $\textup{Var}(X) \leq \frac{(M - m)^2}{4}$.
\end{lemma}

To establish an upper bound on the generalization guarantee of data-driven learning Lagrangian multipliers, we use several tools to control the covering number and Rademacher complexity. For completeness, we provide their formal definitions and statements as follows.
\begin{definition}[Empirical $L^2(\cP^N)$-metric, \citet{wainwright2019high}] \label{def:empirical-l2-pn-metric}
    Let $\cU$ be a function class of which each function takes input in $\cP$. Let $S = \{P_1, \dots, P_N\} \subset \cP$ be a set of $N$ problem instances. We define the \textit{empirical $L^2(\cP^N)$-metric} on $\cU$ as 
    \[
        \|u - u'\|_{2, N} \coloneqq \sqrt{\frac{1}{N}\sum_{i = 1}^N (u(P_i) - u'(P_i))^2}.
    \]
\end{definition}

\begin{definition}[Covering set and covering number, \citet{wainwright2019high}]
    Let $(X, \|\cdot\|)$ be a metric space, and let $\Pi \subset X$ be a subset of $X$. We say that $\Pi' \subset \Pi$ is a $\delta$-covering set of $\Pi$ if for any $\pi \in \Pi$, there exists $\pi' \in \Pi'$ such that $\|\pi - \pi'\| \leq \delta$. A minimum $\delta$-covering set of $\Pi$ is the $\delta$-covering set of $\Pi$ that has the fewest elements. The $\delta$-covering number of $\Pi$, denote $\cN(\delta, \Pi, \|\cdot\|)$ is the number of elements of a minimum $\delta$-covering set. 
\end{definition}

\begin{lemma}[Dudley's entropic integral, \citet{wainwright2019high}] \label{lm:dudley} Let $D = \sup_{u}\|u\|_{2, N}$, we have
    \[
    \hat{\mathscr{R}}_N(\cU) \leq \inf_{\delta_0 > 0} \delta_0 + \int_{\delta_0}^{D} \sqrt{\frac{\log \cN(\delta, \cU, \|\cdot\|_{2, N})}{N}} \mathrm{d}\delta.
    \]
\end{lemma}


\section{Proofs for Section \ref{sec:standard-lagrangian-relaxation}} \label{apx:standard-lagrangian-relaxation}

In this section, we present a detailed proof of the minimax lower bound for the data-driven learning Lagrangian relaxation problem.

\subsection{Additional Backgrounds}


First, we recall a classical result for constructing the packing set for an $s$-dimensional hypercube.

\begin{lemma}[Varshamov-Gilbert, Lemma 2.9, \cite{tsybakov2008nonparametric}] \label{lm:Varshamov-Gilbert-bound}
    Let $s \geq 8$. There exists a subset $\mathcal{M} = \{v^{(1)}, \dots, v^{(M)}\} \subset \{0, 1\}^s$ such that:
    \begin{itemize}
        \item $M \geq 2^{s / 8}$.
        \item $v^{(1)} = (0, \dots, 0) \in \{0, 1\}^s$.
        \item For any distinct pair $v^{(i)}, v^{(j)} \in \mathcal{M}$, the Hamming distance $d_H(v^{(i)}, v^{(j)}) \geq \frac{s}{8}$,
        \item Consequently, the $\ell_1$ distance between any distinct pair satisfies $\|v^{(i)} - v^{(j)}\|_1 \geq \frac{s}{8}$.
    \end{itemize}
\end{lemma}

\subsection{Proof of the Geometric Properties} \label{sec:app-geometric}


\proof[Proof of Proposition~\ref{prop:geometric-property}]
    We first prove the concavity property. Recall that for $P = (c, A, b, C, d)$, the optimal value of the corresponding Lagrangian relaxation $u(\pi, P)$ is defined as
    \[
        u(\pi, P) = \ \min_{x \in \cX} L(\pi, x),
    \]
    where $L(\pi, x) = c^\top x + \pi^\top (b - Ax)$ is the Lagrangian function 
    and $\cX = \{x \in \bbR^m_+ \times \{0, 1\}^p \mid Cx \geq d\}$. Note that $L(\pi, x)$ is an affine function of $\pi$ and therefore also concave in $\pi$. This means that $u(\pi, P)$ is the point-wise minimum of concave functions of $\pi$, hence it is also concave. 
    
    Next, we prove that $g(\pi, P) = b - Ax^*(\pi, P)$ is a valid subgradient of $u(\pi, P)$. First, note that $u(\cdot, P)$ is a concave function of $\pi$. Therefore, a vector $g$ is a valid subgradient of $u$ at $\pi$ if for any multiplier $\pi' \in \bbR^s_+$, we have
    \[
        u(\pi', P) \leq u(\pi, P) + g^\top (\pi' - \pi). 
    \]
    Now let $x^*(\pi, P)$ be the optimal solution that achieves the minimum for $u(\pi, P)$. By definition, we have $u(\pi, P) = c^\top x^*(\pi, P) + \pi^\top (b - Ax^*(\pi, P))$. Now, consider any $\pi' \in \bbR^s_+$, by the definition of $u(\pi', P)$, we have
    \begin{equation*}
        \begin{aligned}
            u(\pi', P) &= \min_{x \in \mathcal X} ~c^\top x + (\pi')^\top (b - Ax) \\
            &\leq c^\top x^*(\pi, P) + (\pi')^\top (b - Ax^*(\pi, P)) \\
            &= \left[c^\top x^*(\pi, P) + \pi^\top (b - Ax^*(\pi, P))\right]  + (\pi' - \pi)^\top (b - Ax^*(x, P)) \\
            &= u(\pi, P) + g(\pi, P)^\top (\pi' - \pi).
        \end{aligned}
    \end{equation*}
    This means that given a problem instance $P$, $g(\pi, P) = b - Ax^*(\pi, P)$ is a valid sub-gradient of $u(\pi, P)$ for any $\pi \in \bbR^s_+$. Finally, from Assumption~\ref{asmp:bounded-constaint-violation}, we have 
    \begin{equation*}
            \|g(\pi, P)\|_2 = \|b - Ax^*(\pi, P)\|_2 \leq \|b\|_2 + \|Ax^*(\pi, P)\|_2 \leq 2B\sqrt{s}.
    \end{equation*}
    The proof is complete.
\qed

\subsection{Minimax Lower-bound for Data-driven Learning Lagrangian Relaxation} \label{sec:app-minimax}

We now present the detailed proof of the minimax lower bound for the data-driven learning Lagrangian relaxation problem. 

\proof[Proof of Theorem~\ref{thm:minimax-lower-bound}]
    We break down the proof into the following steps. 
        
    \textbf{Step 1 - Restricting the problem space:} First, throughout this construction, we consider the problem instance $P$  which is strictly in the form $P = (c, \mathbf{I}_s, \frac{1}{2}\cdot \mathbf{1}_s, 0_{t \times s}, 0_t)$ representing the ILP problem
    \[
        \min_{x \in \{0, 1\}^s} c^\top x \quad \textup{s.t.} \quad x_k \geq \frac{1}{2} \quad \forall k = 1, \dots, s.
    \]
     Here, $(C, d) = (0_{t \times s}, 0_t)$ means all constraints are hard constraints, i.e., they will all be dualized in the Lagrangian relaxation problem. Besides, the problem instance $P$ is solely determined by the objective vector $c$, meaning that we can equate the problem distribution $\cD$ for $P$ as a distribution for $c$. Besides, for any $\pi \in \bbR^s_+$, we have
     \begin{equation*}
         \begin{aligned}
             u(\pi, P) &= \min_{x \in \{0, 1\}^s} c^\top x + \pi^\top (\frac{1}{2}\cdot \mathbf{1}_s - x)  = \frac{1}{2}\pi^\top \mathbf{1}_s + \min_{x \in \{0, 1\}^s}(c - \pi)^\top x  \\
            &= \frac{1}{2}\pi^\top \mathbf{1}_s + \sum_{k = 1}^s \min(0, c_k - \pi_k) = \sum_{k = 1}^s \min \left(\frac{1}{2}\pi_k, c_k - \frac{\pi_k}{2} \right).
         \end{aligned}
     \end{equation*}

    \textbf{Step 2 - Defining problem distribution:} Let $\mu > 0$ be a based constant and $\sigma > 0$ be a scale parameter, where $\mu+\sigma < \pi_{\max}$. For a fixed vector $v \in \{0, 1\}^s$, we define the problem distribution $\cD_v$ over $c$ as follow:
    \begin{itemize}
        \item Support: $c_k$ takes value in $\{\mu, \mu + \sigma\}$ for any $k = 1, \dots, s$,
        \item Probabilities: 
        \begin{itemize}
            \item If $v_k = 1$, then $\bbP(c_k = \mu + \sigma) = \frac{1 + \epsilon}{2} = p_k$, and $\bbP(c_k = \mu) = \frac{1 - \epsilon}{2} = 1 - p_k$, 
            \item If $v_k = 0$, then $\bbP(c_k = \mu + \sigma) = \frac{1 - \epsilon}{2} = p_k$, and $\bbP(c_k = \mu) = \frac{1 + \epsilon}{2} = 1 - p_k$.
        \end{itemize}
    \end{itemize}
    Here, $\epsilon > 0$ is a tunable parameter. Note that with the problem distribution $\cD_v$, the optimal Lagrangian relaxation $\pi^*(\cD_v)$ is
    \begin{equation*}
        \begin{aligned}
            \pi^*(\cD_v) = \arg \max_{\pi \in \Pi} \bbE_{c \sim \cD_v} \Big[\sum_{k = 1}^s \min\left(\frac{1}{2}\pi_k, c_k - \frac{\pi_k}{2}\right) \Big].
        \end{aligned}
    \end{equation*}
    Note that this optimization problem can be decomposed into $s$ sub-problems, i.e., $(\pi^*(\cD_v))_k$ can be defined as
    \begin{equation*}
        \begin{aligned}
            (\pi^*(\cD_v))_k &= \arg \max_{\pi_k} \left\{ \mathcal J_k(\pi_k) \triangleq \bbE_{c_k \sim \cD_{v_k}} \Big[ \min\left(\frac{1}{2}\pi_k, c_k - \frac{\pi_k}{2}\right)\Big] \right\}
        \end{aligned}
    \end{equation*}
    We then have the following cases
    \begin{itemize}
        \item Case 1: $v_k = 1$, then $p_k > \frac{1}{2}$. We have
        \begin{equation*}
            \begin{aligned}
                \mathcal J_k(\pi_k) &= p_k \cdot\min\left(\frac{1}{2}\pi_k, \mu + \sigma - \frac{\pi_k}{2}\right) + (1 - p_k)\cdot \min\left(\frac{1}{2}\pi_k, \mu - \frac{\pi_k}{2}\right),
            \end{aligned}
        \end{equation*}
        and we consider the following cases:
        \begin{itemize}
            \item Case 1.1: $\pi_k < \mu$, then $\mathcal J_k(\pi_k) = \frac{\pi_k}{2}$ and $\frac{\partial \mathcal J_k}{\partial \pi_k} = \frac{1}{2}$. This means $\mathcal J_k$ is increasing over $[0, \mu]$.
            \item Case 1.2: $\pi_k > \mu + \sigma$, then $\mathcal J_k(\pi_k) = p_k\left(\mu + \sigma - \frac{\pi_k}{2}\right) + (1 - p_k)\left(\mu -  \frac{\pi_k}{2}\right)$, and $\frac{\partial \mathcal J_k}{\partial \pi_k} = -\frac{1}{2}$. This means $\mathcal J_k$ is decreasing over $[\mu + \sigma, \infty)$.
            \item Case 1.3: $\pi_k \in [\mu, \mu + \sigma]$, then $\mathcal J_k(\pi_k) = p_k\cdot \frac{\pi_k}{2} + (1 - p_k)(\mu - \frac{\pi_k}{2})$ and $\frac{\partial \mathcal J_k}{\partial \pi_k} = p_k - \frac{1}{2}$. Since $p_k > \frac{1}{2}$, this means $\mathcal J_k$ is increasing over $[\mu, \mu + \sigma]$.
        \end{itemize}
        Therefore, $\mathcal J_k$ attains its maxima at $\pi_k = \mu + \sigma$, i.e., $(\pi^*(\cD_v))_k = \mu + \sigma$.

        \item Case 2: $v_k = 0$, then $p_k < \frac{1}{2}$. Using the same idea as the previous case, we have $(\pi^*(\cD_v))_k = \mu$.
    \end{itemize}
    Then, we conclude that 
    \begin{equation*}
        \begin{aligned}
            \pi^*(\cD_v) &= \arg\max_{\pi \in \Pi}\bbE_{c \sim \cD_v} \sum_{k = 1}^s \min\left(\frac{1}{2}\pi_k, c_k - \frac{\pi_k}{2}\right) = \mu \mathbf{1}_s + \sigma v.
        \end{aligned}
    \end{equation*}

    \textbf{Step 3 - Defining a family of hard problem distributions}: Let $\{0, 1\}^s$ be an $s$-dimensional hypercube. From Lemma \ref{lm:Varshamov-Gilbert-bound}, there exists a subset $\cV = \{v^{(1)}, \dots, v^{(M)}\} \subset \{0, 1\}^s$ such that $M \geq 2^{s/8}$ and $\|v^{(i)} - v^{(j)}\|_1 \geq \frac{s}{8}$. This means that for any distinct $v^{(i)}, v^{(j)} \in \cV$, we have
    \begin{equation*}
        \begin{aligned}
            \|\pi^*(\cD_{v^{(i)}}) - \pi^*(\cD_{v^{(j)}})\|_1 &= \|\mu \mathbf{1}_s + \sigma v^{(i)} - \mu \mathbf{1}_s - \sigma v^{(j)}\|_1 = \sigma \|v^{(i)} - v^{(j)} \|_1 \geq \frac{\sigma s}{8}.
        \end{aligned}
    \end{equation*}
    This means that the set $\{\pi^*(\cD_{v^{(1)}}), \dots, \pi^*(\cD_{v^{(M)}})\}$ is a $\frac{\sigma s}{16}$-packing set of $\Pi$. Besides, the set $\cV$ defines a family of hard problem distribution $\{\cD_{v^{(1)}}, \dots, \cD_{v^{(M)}}\}$, and the minimax risk is lower-bounded by
    \begin{equation*}
        \begin{aligned}
            &\inf_{\pi}\sup_{\cD \in \Delta(\mathcal{P})} \bbE_{S \sim \cD^N}[U(\pi^*(\cD), \cD) - U(\pi(S), \cD)] 
            \geq &\inf_{\pi} \sup_{v \in \cV} \bbE_{S \sim \cD_v^N}[U(\pi^*(\cD_v), \cD_v) - U(\pi(S), \cD_v)].
        \end{aligned}
    \end{equation*}


    \textbf{Step 4 - Linking the risk with the distance in the parameter space}: We now show that for any $\pi_k$, we have
    \[
        \mathcal J_k((\pi^*(\cD_v))_k) - \mathcal J_k(\pi_k) \geq \frac{\epsilon}{2}\abs{(\pi^*(\cD_v))_k - \pi_k}.
    \]
    If $v_k = 1$, we consider the following case:
    \begin{itemize}
        \item If $\pi_k < \mu$: We have
        \begin{equation*}
            \begin{aligned}
                \mathcal J_k((\pi^*(\cD_v))_k) - \mathcal J_k(\pi_k) 
                = & \mathcal J_k((\pi^*(\cD_v))_k) - \mathcal J_k(\mu) + \mathcal J_k(\mu) - \mathcal J_k(\pi_k) \\
                = &\frac{\epsilon}{2}((\pi^*(\cD_v))_k - \mu) + \frac{1}{2}(\mu - \pi_k) \\
                \geq &\frac{\epsilon}{2}((\pi^*(\cD_v))_k - \mu) + \frac{\epsilon}{2}(\mu - \pi_k) = \frac{\epsilon}{2}\left((\pi^*(\cD_v))_k - \pi_k\right).
            \end{aligned}
        \end{equation*}
        \item If $\pi_k > \mu + \sigma$: We have
        \begin{equation*}
            \begin{aligned}
                \mathcal J_k((\pi^*(\cD_v))_k) - \mathcal J_k(\pi_k)
                = &p_k\cdot \frac{\mu + \sigma}{2} + (1 - p_k)\left(\mu - \frac{\mu + \sigma}{2}\right) 
                - p_k\left(\mu + \sigma - \frac{\pi_k}{2}\right) - (1 - p_k)\left(\mu - \frac{\pi_k}{2}\right) \\
                = &p_k\left(- \frac{\mu + \sigma}{2} + \frac{\pi_k}{2}\right) + (1 - p_k) \left(\frac{\pi_k}{2} - \frac{\mu + \sigma}{2}\right) \\
                = &\frac{\pi_k}{2} - \frac{\mu + \sigma}{2} \\
                = &\frac{1}{2}(\pi_k - (\pi^*(\cD_v))_k) \qquad  \qquad
                 (\textup{as } (\pi^*(\cD_v))_k = \mu + \sigma) \\
                \geq &\frac{\epsilon}{2}(\pi_k - (\pi^*(\cD_v))_k) \qquad  \qquad (\textup{as $\epsilon$ is small}).
            \end{aligned}
        \end{equation*}
        \item If $\pi_k \in [\mu, \mu + \sigma]$: From the above, in such a case, we have $\frac{\partial \mathcal J_k}{\partial \pi_k} = p_k - \frac{1}{2} = \frac{\epsilon}{2}$ for any $\pi_k \in [\mu, \sigma]$. And since $\mathcal J_k$ is linear over $[\mu, \mu + \sigma]$ and $(\pi^*(\cD_v))_k = \mu + \sigma$, we claim that $\mathcal J_k((\pi^*(\cD_v))_k) - \mathcal J(\pi_k) = \frac{\epsilon}{2}\abs{(\pi^*(\cD_v))_k - \pi_k}$.
    \end{itemize}
    Therefore, if $v_k = 1$, we can claim that $\mathcal J_k((\pi^*(\cD_v))_k) - \mathcal J_k(\pi_k) \geq \frac{\epsilon}{2}\abs{(\pi^*(\cD_v))_k - \pi_k}$. Using an analogous argument, we can claim that for $v_k = 0$, we also have $\mathcal J_k((\pi^*(\cD_v))_k) - \mathcal J_k(\pi_k) \geq \frac{\epsilon}{2}\abs{(\pi^*(\cD_v))_k - \pi_k}$. Take the sum across the index $k = 1, \dots, s$, and note that $P$ is determined by only the objective vector $c$ (hence we can write $c \sim \cD_v$ instead of $P \sim \cD_v$), we have
    \begin{equation*}
        \begin{aligned}
        U(\pi^*(\cD_v), \cD_v) - U(\hat{\pi}, \cD_v)
            = &\bbE_{c \sim \cD_v}\left[u(\pi^*(\cD_v)) - u(\pi, P)\right] \\
            = &\sum_{k = 1}^s \bbE_{c \sim \cD_v} [\mathcal J_k((\pi^*(\cD_v))_k) - \mathcal J_k(\pi_k)] \ge \frac{\epsilon}{2}\|(\pi^*(\cD_v) - \pi\|_1.
        \end{aligned}
    \end{equation*}
    Combining with the claim from Step 3, we have
    \begin{equation*}
        \begin{aligned}
            \inf_{\pi}\sup_{\cD \in \Delta(\mathcal{P})} \bbE_{S \sim \cD^N}[U(\pi^*(\cD), \cD) - U(\pi(S), \cD)]
            \geq \frac{\epsilon}{2}\inf_{\pi} \sup_{v \in \cV} \bbE_{S \sim \cD_v^N}\|\pi^*(\cD_v) - \pi(S)\|_1.
        \end{aligned}
    \end{equation*}
    
    \textbf{Step 5: Applying Fano's method (Theorem \ref{thm:fano}):} We now use Fano's method to give a lower bound for the LHS above. We have 
    \begin{itemize}
        \item From Lemma \ref{lm:Varshamov-Gilbert-bound}, the log cardinality of the packing set $\mathcal V$ is lower-bounded by $\log M \geq \frac{s}{8}\log 2$.
        \item \textit{Mutual information upper-bound}: In our construction, the coordinate of objective vector $c$ are independent, the KL divergence sums over $s$ dimensions. For any distinct $v^{(i)}, v^{(j)} \in \mathcal V$, the KL divergence between $\cD_{v^{(i)}}$ and $\cD_{v^{(j)}}$ is upper-bounded by the sum of $s$ Bernoulli KL divergences. {For small $\epsilon \in (0, 1/2)$, using standard $\chi^2$-upper bound for KL divergence, we have $\textup{KL}(\textup{Ber}(\frac{1 + \epsilon}{2}) \| \textup{Ber}(\frac{1 - \epsilon}{2})) \leq \chi^2(P \| Q) =  \frac{4\epsilon^2}{1 - \epsilon^2} = \cO(\epsilon^2)$ using $\chi^2$-upper bound.} Therefore, we have $I(J; S) \leq N \max_{i, j} \textup{KL}(\cD_{v^{(i)}} \| \cD_{v^{(j)}}) \leq  N s (4\epsilon^2)$.
    \end{itemize}
    
    To make Fano's inequality valid, we need to choose $\epsilon$ such that
    \[
        \frac{4Ns \epsilon^2 + \log 2}{\frac{s}{8}\log 2} \leq \frac{1}{2} \Leftrightarrow 4Ns\epsilon^2 \leq (\frac{s}{16} - 1)\log 2.
    \]
    We simply choose $\epsilon \asymp \frac{1}{\sqrt{N}}$. Combining with the claim from Steps 3, 4, and applying Theorem~\ref{thm:fano}, we have
    \begin{equation*}
        \begin{aligned}
            \inf_{\pi}\sup_{\cD \in \Delta(\mathcal{P})} \bbE_{S \sim \cD^N}[U(\pi^*(\cD), \cD) - U(\pi(S), \cD)]
            \geq &\frac{\epsilon}{2}\inf_{\pi} \sup_{v \in \cV} \bbE_{S \sim \cD_v^N}\|\pi^*(\cD_v) - \pi(S)\|_1 = \Omega \left(\frac{1}{\sqrt{N}} \cdot \sigma s\right),
        \end{aligned}
    \end{equation*}
    which concludes the proof.
\qed

\subsection{Minimax Optimality via Stochastic Gradient Ascent} \label{apx:SGA}
In this section, we will present the formal proof for Theorem \ref{thm:minimax-SGA}.

\begin{customthm}{\ref{thm:minimax-SGA} (restated)}
    If the SGA algorithm (Algorithm \ref{alg:sGA}) is run for $N$ iterations with a constant step size $\eta = \frac{\pi_{\max}}{2B\sqrt{N}}$, the expected excess risk of the learned (averaged) Lagrangian multipliers $\bar{\pi}_N$ is upper-bounded by
    \[
        \bbE_{S \sim \cD^N}[\mathcal{E}(\bar{\pi}_N)] \leq \frac{2B\pi_{\max} s}{\sqrt{N}} = \cO\left(\frac{s}{\sqrt{N}}\right).
    \]
\end{customthm}
\proof[Proof of Theorem~\ref{thm:minimax-SGA}]
    Let $U(\pi) = \bbE_{P \sim \cD}[u(\pi, P)]$ denote the expected utility function corresponding to Lagrangian variables $\pi$. Since $u(\cdot, P)$ is concave for any problem instance $P$, $U(\pi)$ is also concave. Besides, recall that $\pi^*(\cD) \in \arg \max_{\pi \in \Pi} U(\pi)$ is the optimal Lagrangian variable  corresponding to the problem distribution $\cD$.

    From the property of the projection operator onto a convex set $\Pi$, for any $t  \in \{1, \dots, N\}$, we have
    \begin{equation} \label{eq:descent-lemma}
        \begin{aligned}
            &\|\pi_{t + 1} - \pi^*(\cD)\|_2^2 \leq \|\pi_{t} + \eta g_t - \pi^*(\cD)\|_2^2 = \|\pi_t - \pi^*(\cD)\|_2^2 + 2\eta g_t^\top (\pi_t - \pi^*(\cD)) +\eta^2 \|g_t\|_2^2 \\
            \Rightarrow & \quad g_t^\top (\pi^*(\cD) - \pi_t) \leq \frac{\|\pi_t - \pi^*(\cD)\|_2^2 - \|\pi_{t + 1} - \pi^*(\cD)\|_2^2}{2\eta} + \frac{\eta}{2} \|g_t\|_2^2.
        \end{aligned}
    \end{equation}
    Note that $u(\cdot, P_t)$ is a concave function, therefore $u(\pi^*(\cD), P_t) - u(\pi_t, P_t) \leq g^\top_t (\pi^*(\cD) - \pi_t)$. Note that $\pi_t$ is constructed using only the historical instances $\{P_1, \dots, P_{t - 1}\}$, hence is independent with $P_{t}$. Therefore, if we take the expectation over all problem instances $P_1, \dots, P_N \sim \cD$, then $\bbE_{P_1,\dots, P_N}[u(\pi_t, P)] = \bbE_{P_1,\dots, P_N}[U(\pi_t)]$, and $\bbE_{P_1,\dots, P_N}[u(\pi^*(\cD), P)] = U(\pi^*)$ since $\pi^*$ is deterministic. Combining the above with Equation \ref{eq:descent-lemma}, we have
    \begin{equation*}
        \bbE_{P_1,\dots, P_N}[U(\pi^*) - U(\pi_t)] \leq \bbE_{P_1,\dots, P_N}\left[\frac{\|\pi_t - \pi^*(\cD)\|_2^2 - \|\pi_{t+ 1} - \pi^*(\cD)\|_2^2}{2\eta} + \frac{\eta}{2}\|g_t\|_2^2\right].
    \end{equation*}
    Taking the sum over $N$ time steps, we have
    \begin{equation*}
        \begin{aligned}
            \bbE_{P_1,\dots, P_N}\left[N U(\pi^*) - \sum_{t = 1}^N U(\pi_t)\right] &\leq  \bbE_{P_1,\dots, P_N} \left[\frac{\|\pi_1 - \pi^*(\cD)\|_2^2 - \|\pi_{N + 1} - \pi^*(\cD)\|_2^2}{2\eta} + \frac{\eta}{2}\sum_{t = 1}^N \|g_t\|_2^2\right] \\
            &\leq  \frac{\|\pi_1 - \pi^*(\cD)\|_2^2}{2\eta} + \bbE_{P_1,\dots, P_N}\left[\sum_{t = 1}^N \|g_t\|_2^2\right],
        \end{aligned}
    \end{equation*}
    where we use the fact that $\pi_1$ is fixed, so $\bbE_{P_1,\dots, P_N}\|\pi_1 - \pi^*(\cD)\|_2^2 = \|\pi_1 - \pi^*(\cD)\|_2^2$. Now note that $\|g_t\|_2 \leq 2B\sqrt{s}$, and $\|\pi_1 - \pi^*(\cD)\|_2^2 \leq s\pi_{\max}^2$, we have
    \[
        \bbE_{P_1,\dots, P_N}\left[U(\pi^*) - \sum_{t = 1}^N U(\pi_t)\right] \leq \frac{s\pi_{\max}^2}{2\eta} + 2\eta NB^2s.
    \]
    Now, since $U(\cdot)$ is a concave function, using Jensen's inequality, we have $\frac{1}{N}\sum_{t = 1}^N U(\pi_t) \leq U\left(\frac{1}{N}\sum_{t =1}^N \pi_t\right) = U(\overline{\pi}_N)$, and therefore
    \[
       N \bbE[\cE(\bar{\pi}_N)] =  N \bbE_{P_1, \dots, P_N}[U(\pi^*(\cD)) - U(\bar{\pi}_N)] \leq \frac{s\pi_{\max}^2}{2\eta} + 2\eta NB^2s.
    \]
    Choosing $\eta = \frac{\pi_{\max}}{2B\sqrt{N}}$, we have the final conclusion. 
\qed

\section{Proofs for Section \ref{sec:warm-start-lagrangian-relaxation}} \label{apx:warm-start-lagrangian-relaxation}
\subsection{Upper Bound} \label{apx:upper-bound-warm-start}
We now present the detailed proof for the generalization guarantee upper-bound for the problem of learning to warm-start Lagrangian relaxation. 

\proof[Proof of Theorem~\ref{thm:upper-bound-warm-start}]
    The proof simply relies on the closed-form solution of the ERM estimator for the squared loss. First, given an initialization $\phi \in \Pi$, the risk $R(\phi) = \bbE_{P \sim \cD}[\|\phi - \pi^*(P)\|_2^2]$. Therefore, given a problem distribution $\cD$, the optimal initialization $\phi^*(\cD)$ corresponding to $\cD$ is simply $\phi^*(\cD) = \bbE_{P \sim \cD}[\pi^*(P)]$. Therefore, the excessive can be written as 
    \[
        \bbE_{P \sim \cD}[\|\hat{\phi}(S) - \pi^*(P)\|_2^2] - \bbE_{P \sim \cD}[\|\pi^*(\cD) - \pi^*(P)\|] = \|\hat{\phi}(S) - \phi^*\|_2^2.
    \]
    Let $Z_i = \pi^*(P_i)$, and from Assumption \ref{asmp:restricted-search-main}, the $s$-dimensional random vector $Z_i$ are i.i.d.~with mean $\phi^*(\cD)$ and belongs to the bounded domain $Z_i \in [0, \pi_\textup{max}]^s$ due to Assumption \ref{asmp:restricted-search-main}. { Therefore, applying the standard Popoviciu's inequality on variance (Lemma \ref{lm:popoviciu}), we have}
    \[
        \bbE_{P \sim \cD} = \frac{1}{N}\sum_{i = 1}^N \sum_{k = 1}^s \textup{Var}(Z_{i, k})\leq \frac{s \pi^2_\textup{max}}{4N} = \cO\left(\frac{s}{N}\right),
    \]
    where $Z_{i, k}$ is the $k$-th coordinate of $Z_i$.
\qed

\subsection{Lower Bound}
We now present a detailed proof of the minimax lower bound for the problem of learning to warm-start Lagrangian relaxation. 

\begin{customthm}{\ref{thm:minimax-lower-bound-warm-start} (restated)}
    For any learning algorithm $\hat{\phi}$ taking input as a set $S$ of $N$ problem instances $P_1, \dots, P_N \sim \cD$, the minimax expected risk is lower-bounded by
    \[
        \inf_{\phi} \sup_{\cD \in \Delta(\cP)}\bbE_{S \sim \cD^N}\left[ \bbE_{P \sim \cD}[\| \phi(S) - \pi^*(P)\|_2^2] - \bbE_{P \sim \cD}[\|\phi^* - \pi^*(P)\|_2^2] \right] = \Omega\left(\frac{s}{N}\right). 
    \]
\end{customthm}
\begin{proof}[Proof of Theorem~\ref{thm:minimax-lower-bound-warm-start}]
    Similar to Theorem \ref{thm:minimax-lower-bound}, our approach is to leverage Fano's method (Theorem \ref{thm:fano}). However, since the loss here is quadratic, the overall proof can be simplified drastically. We proceed with the following steps.

    \textbf{Step 1: Construction of hard instances.} In this construction, we also leverage the idea from the previous case (Theorem \ref{thm:minimax-lower-bound}), by restricting the problem instance to take only the form $P = (c, \mathbf{I}_s, \frac{1}{2}\cdot \mathbf{1}_s, 0_{t \times s}, 0_t)$. From Lemma \ref{lm:lower-bound-step-1}, we show that problem instance $P$ of this form has the property $u(\pi, P) = \sum_{k = 1}^s \min(\frac{\pi_k}{2}, c_k - \frac{\pi_k}{2})$. This function is strictly concave and is maximized uniquely when $\frac{\pi_k}{2} = c_k - \frac{\pi_k}{2}$, if $c_k > 0$ since the Lagrangian multiplier has to be positive. Therefore, for any $c$ such that $c_k \geq 0$ for all $k$, we have $\pi^*(P) = c$.

     \textbf{Step 2: Family of distribution construction.} We define a family of distributions for $c$, parameterized by a binary vector $v \in \{0, 1\}^s$ and a small perturbation scale $\epsilon \in (0, 1/2)$ that we will choose later.
     \begin{itemize}
         \item Support: for any $k$, $c_k$ takes value in $\{1, 2\}$.
         \item Probability: the probabilities are controlled by $v$ and $\epsilon$ as follows
         \begin{itemize}
             \item If $v_k = 1$: $\bbP(c_m = 2) = \frac{1 + \epsilon}{2}$ and $\bbP(c_k = 1) = \frac{1 - \epsilon}{2}$.
             \item If $v_k = 0$: $\bbP(c_k = 2) = \frac{1 + \epsilon}{2}$ and $\bbP(c_k = 1) = \frac{1 + \epsilon}{2}$.
         \end{itemize}
     \end{itemize}
     Then, for a problem distribution $\cD_{v}$ constructed this way, we have $(\phi^*(\cD_v))_k = \bbE[c_k] = \frac{3}{2} + \frac{\epsilon}{2}(2v_k - 1)$
     
     \textbf{Step 4: Minimax lower-bound.} From Varshamov-Gilbert bound in Lemma~\ref{lm:Varshamov-Gilbert-bound}, there is a packing set $\cV = \{v^{(1)}, \dots, v^{(M)}\} \subset \{0, 1\}^s$ with $M \geq 2^{s / 8}$ such that $d_H(v^{(i)}, v^{(j)}) \geq s / 8$. Therefore, for any two $v^{(i)}, v^{(j)} \in \cV$, we have 
     \begin{equation} \label{eq:lower-bound-warm-start-critical-ineq}
         \|\phi^*(\cD_{v^{(j)}}) - \phi^*(\cD_{v^{(i)}})\|_2^2 = \sum_{k = 1}^s\left(\frac{\epsilon}{2}(2v_k^{(j)}) - 2v_{k}^{(i)})\right)^2 = \epsilon^2 \cdot d_H(v^{(j)}, v^{(i)}) \geq \frac{\epsilon^2 s}{8}.
     \end{equation}

     Again, we note that the KL divergence for the Bernoulli distributions with parameter $p, q \in [\frac{1 - \epsilon}{2}, \frac{1 + \epsilon}{2}]$ satisfies $\textup{KL}(p \parallel q) \leq \frac{(p - q)^2}{q(1 - q)} \leq 4 \epsilon^2$. Therefore,
     \[
        \textup{KL}(\cD_{v^{(i)}}^N \parallel\cD_{v^{(j)}}^N) \leq 4N\sum_{k = 1}^s (p_k^{(j)} - p_k^{(i)})^2 = 4N\epsilon^2 d_H(v^{(i)}, v^{(j)}) \leq 4Ns\epsilon^2.
     \]
     We choose $\epsilon$ such that
     \[
        4Ns\epsilon^2 \leq \frac{s}{16}\log 2 \Rightarrow \epsilon \asymp \frac{1}{\sqrt{N}}.
     \]
     Substituting to~\eqref{eq:lower-bound-warm-start-critical-ineq}, we conclude that 
     \[
        \inf_{\phi} \sup_{\cD \in \Delta(\cP)}\bbE_{S \sim \cD^N}\left[ \bbE_{P \sim \cD}[\| \phi(S) - \pi^*(P)\|_2^2] - \bbE_{P \sim \cD}[\|\phi^* - \pi^*(P)\|_2^2] \right] = \Omega\left(\frac{s}{N}\right).
     \]

     The proof is complete.
\end{proof}



